\section{结论}

任务信任将关于物理观测、受保护记录、共同治理状态，以及特定用途适合性的证据联系起来。OAL 按对象和生命周期组织这些主张，带类型的依赖图与 CTG 循环则识别何时需要因变化而重新评估。将框架应用于 39 项全文研究，区分了互补机制、各自的评估环境，以及每次交接所需的证据。分层与敏感性分析使覆盖模式对选择和编码的依赖变得清晰。

分区-撤销案例明确了联盟角色、面向 Fabric 的应用模式、受限验证和可复现通信预算。其可执行回放保留了本地权限、事件时刻撤销和后续审计之间的区别。具有策略验证能力的见证日志在匹配规则下产生相同决策。因此，账本价值取决于共享治理需求，以及实现这些需求的实测成本。

由此形成的研究议程要求独立可复现的编码、匹配的架构与评估基线、真实的分区与恢复测试，以及集成任务证据。该框架为将机制证据与任务要求相匹配，并选择这些要求所支持的治理服务，提供可追溯的基础。
